% ===================================================================
%  TABLICA Nr 4 — Wiek dojrzały i starość.
%  Meme regime que les precedents.
%
%  LA PARENTE PAR ALLIANCE EST UN SYSTEME, ET LE POLONAIS L'A ENTIER :
%  teść, teściowa, zięć, synowa, szwagier, bratowa, szwagierka. L'ido
%  la fabrique avec « bo- » — bopatrulo, bofilio, bofratino — et le
%  francais avec « beau- ». Chacun de ces mots porte un renvoi ; les
%  confondre, c'est faire pointer la mauvaise personne.
%
%  ET « wnuk » N'EST PAS « wnuczek ». Le releve distingue nepotuli et
%  nepotino ; le polonais dit wnuk / wnuczka, et le diminutif
%  « wnuczek » n'est pas le feminin — c'est le petit-fils qu'on
%  cherit. On ne l'emploie donc pas ici.
% ===================================================================

%%K t04-apar-1 apar
\VUsaut{4.0mm}
\VUtitre{15.0pt}{60}{TABLICA Nr 4}
\VUsaut{1.6mm}
\VUtitre{11.4pt}{0}{Wiek dojrzały i starość.}
\VUsaut{3.2mm}

%%K t04-tit-1 sub
\VUcentre{11.4pt}{0}{\textit{Scena pierwsza.}}
\VUsaut{1.6mm}
\VUtitre{11.4pt}{0}{Uczta weselna.}
\VUsaut{2.6mm}

%%K t04-tit-2 sub
\VUpk{10.2pt}{40}{Jesień.}
\VUsaut{3.0mm}

%%K t04-c1-01-1 p
1. --- Młodzi urośli. Jeden z nich, Jan, się żeni. Jesteśmy na \VUgras{uczcie
weselnej}, która gromadzi wokół tego samego stołu rodziny małżonków.

%%K t04-c1-02-1 p
2. --- Jesteśmy w \VUgras{jesieni}, w miesiącu \VUgras{wrześniu}. Zimny wiatr
\VUgras{października} i \VUgras{listopada} nie ogołocił jeszcze wsi z jej zielonego
listowia. Wieczorne powietrze jest ciepłe i wonne; dlatego obiad odbywa się pod
\VUgras{werandą} \textsuperscript{(8)} przy ogrodzie.

%%K t04-c1-03-1 p
3. --- Daleko, na \VUgras{wzgórzu} \textsuperscript{(3)}, znajduje się dom rodziców
Jana, elegancki \VUgras{zamek} \textsuperscript{(1)} ukryty w zieleni; piękne winnice,
które dają wyborne wino, okrywają stok pagórka. Winiarz je uprawia i pielęgnuje.

%%K t04-c1-04-1 p
4. --- Nad winnicami przelatuje stado \VUgras{kuropatw} \textsuperscript{(2)}
przestraszonych zbliżaniem się \VUgras{gajowego} \textsuperscript{(5)}. Ten, ze
\VUgras{strzelbą} pod pachą i z \VUgras{torbą myśliwską} przez ramię, przeszukuje
\VUgras{zarośla} \textsuperscript{(4)} ze swoim \VUgras{psem} \textsuperscript{(6)}.
Dogląda majątku i przeszkadza kłusownikom niszczyć zwierzynę.

%%K t04-c1-05-1 p
5. --- Na zewnątrz \VUgras{werandy} pnie się \VUgras{winorośl}, z której zwisają gęste
\VUgras{grona} \textsuperscript{(7)}.

%%K t04-c1-06-1 p
6. --- Piękne kwiaty jesienne, które zdobią \VUgras{stół} \textsuperscript{(16)}, to
\VUgras{chryzantemy} \textsuperscript{(9)}; \VUgras{ojciec} \textsuperscript{(10)} panny
młodej wpiął jedną z nich w dziurkę swojego fraka.

%%K t04-c1-07-1 p
7. --- Posiłek się kończy. \VUgras{Służący} \textsuperscript{(11)} zaczyna przynosić
półmiski z deserem i wkrótce zabierze różne części zastawy, \VUgras{łyżki}
\textsuperscript{(14)}, \VUgras{widelce} \textsuperscript{(12)}, \VUgras{noże}
\textsuperscript{(13)}, \VUgras{półmiski} \textsuperscript{(15)}, których już nie
potrzeba.

%%K t04-tit-3 sub
\VUpk{10.2pt}{40}{Krewni.}
\VUsaut{3.0mm}

%%K t04-c1-08-1 p
8. --- Wszyscy wyglądają radośnie, a uśmiech, z jakim \VUgras{matka}
\textsuperscript{(17)} pana młodego patrzy na swoje dzieci, dowodzi jej szczęścia.

%%K t04-c1-09-1 p
9. --- To małżeństwo łączy dwie bogate rodziny z okolicy. Jan, \VUgras{mąż}
\textsuperscript{(18)}, jest \VUgras{synem} wielkiego właściciela ziemskiego i staje się
\VUgras{zięciem} zdolnego przemysłowca.

%%K t04-c1-10-1 p
10. --- \VUgras{Małżonkowie} są młodzi, a jednak od dawna byli \VUgras{zaręczeni}.
\VUgras{Młoda kobieta} \textsuperscript{(19)}, Marta, jest czarująca w swojej
\VUgras{białej sukni} przybranej kwiatem pomarańczy.

%%K t04-c1-11-1 p
11. --- Jej \VUgras{teść} \textsuperscript{(20)} jest bardzo szczęśliwy, że ma tak miłą
\VUgras{synową}, a \VUgras{teściowa} \textsuperscript{(21)} Jana uśmiecha się, widząc
swoją \VUgras{córkę} tak piękną.

%%K t04-c1-12-1 p
12. --- Na końcu stołu siedzą pozostali \VUgras{goście} \textsuperscript{(22)},
\VUgras{krewni, przyjaciele, kuzyni, kuzynki}. \VUgras{Marszałek stołu}
\textsuperscript{(23)}, z serwetą pod pachą, gorliwie ich obsługuje.

%%K t04-tit-4 sub
\VUpk{10.2pt}{40}{Przyjaciele.}
\VUsaut{3.0mm}

%%K t04-c1-13-1 p
13. --- Po prawej stronie stołu oto \VUgras{panna} \textsuperscript{(24)},
\VUgras{przyjaciółka} młodej żony, potem \VUgras{brat} tejże, który jest jej
\VUgras{drużbą} \textsuperscript{(25)}. Rozmawia ze swoją \VUgras{druhną}
\textsuperscript{(26)}, siostrą męża, która przez to małżeństwo staje się
\VUgras{bratową} Marty. Jest czarującą młodą osobą. Ma piękne \VUgras{klejnoty,
naszyjnik z pereł} i złotą \VUgras{bransoletkę}.

%%K t04-c1-14-1 p
14. --- Obok nich pan, który stoi i podnosi swój kieliszek, jest \VUgras{przyjacielem}
\textsuperscript{(27)} rodziny. Jest jednym ze \VUgras{świadków pana młodego}. Wznosi
\VUgras{toast}, pije za \VUgras{zdrowie} młodych małżonków i życzy im długiego
szczęścia. Jest ubrany w strój uroczysty, we frak, z \VUgras{lakierowanymi butami}
\textsuperscript{(28)}. Towarzyszy \VUgras{babce} \textsuperscript{(29)}, która jest
\VUgras{wdową}.

%%K t04-c1-15-1 p
15. --- Młodzieniec we \VUgras{fraku} \textsuperscript{(31)}, z cienkim czarnym wąsem,
jest \VUgras{szwagrem} \textsuperscript{(30)} Jana. Jest wybitnym inżynierem. Siedzi
obok \VUgras{młodej pani} \textsuperscript{(33)} bardzo eleganckiej, \VUgras{starszej
siostry} Marty i matki ślicznej dziewczynki, która przychodzi, podając ramię swojemu
kuzynowi, ofiarować kwiat i \VUgras{życzyć} jej \VUgras{dobrego wieczoru}. Ta pani ma
bardzo piękne \VUgras{kolczyki} i elegancki \VUgras{strój głowy}.

%%K t04-c1-16-1 p
16. --- Mały kuzyn, tak dziwny w swojej \VUgras{bluzie} \textsuperscript{(36)} i w
swoich bufiastych \VUgras{spodenkach} \textsuperscript{(37)}, jest bardzo godny
politowania. Jest \VUgras{sierotą} \textsuperscript{(35)}. Bardzo młodo stracił ojca i
matkę. Jego wuj jest jego \VUgras{opiekunem} \textsuperscript{(38)} i został kawalerem,
żeby wychować biedne dziecko. Jest dobrym człowiekiem. Jest trochę łysy i bardzo
krótkowzroczny; używa \VUgras{binokli}.

%%K t04-tit-5 sub
\VUsaut{2.4mm}
\VUpk{10.2pt}{40}{Deser.}
\VUsaut{3.0mm}

%%K t04-c1-17-1 p
17. --- Za biesiadnikami, na stole nakrytym \VUgras{obrusem}, przygotowany jest
\VUgras{deser} \textsuperscript{(39)}. W wielkiej czarze oto owoce, \VUgras{brzoskwinie}
\textsuperscript{(40)}, \VUgras{gruszki} \textsuperscript{(41)}, \VUgras{jabłka}
\textsuperscript{(42)}, \VUgras{morele} \textsuperscript{(43)} i \VUgras{winogrona}
\textsuperscript{(44)}. Obok \VUgras{ananasy} \textsuperscript{(45)}, piękne
\VUgras{ciasto} \textsuperscript{(46)}, \VUgras{butelki likieru} \textsuperscript{(48)}
i \VUgras{szampan} \textsuperscript{(49)}, a także stosy czystych \VUgras{talerzy}
\textsuperscript{(47)}.

%%K t04-c1-18-1 p
18. --- \VUgras{Piwniczny} \textsuperscript{(50)} odkorkowuje \VUgras{butelkę}
\textsuperscript{(52)} starego wina \VUgras{korkociągiem} \textsuperscript{(51)}.
\VUgras{Korek} \textsuperscript{(53)} wydaje się bardzo trudny do wyciągnięcia. Ten
piwniczny ma \VUgras{serwetę} \textsuperscript{(54)} pod pachą.

%%K t04-c1-19-1 p
19. --- Po obiedzie panowie pójdą do palarni, a panie pójdą się przygotować na bal.

%%K t04-tit-6 sub
\VUsaut{5.0mm}
\VUcentre{11.4pt}{0}{\textit{Scena druga.}}
\VUsaut{1.6mm}
\VUpk{11.4pt}{40}{Urodziny dziadka.}
\VUsaut{2.6mm}

%%K t04-tit-7 sub
\VUpk{10.2pt}{40}{Odwiedziny.}
\VUsaut{3.0mm}

%%K t04-c2-01-1 p
1. --- Minęło dziesięć lat, rodzice Jana zestarzeli się, bardzo kochając swoje troje
wnucząt, dwóch \VUgras{wnuków} \textsuperscript{(2)} i jedną \VUgras{wnuczkę}
\textsuperscript{(3)}. Ponieważ dzisiaj przypadają \VUgras{urodziny dziadka}, dzieci
przychodzą złożyć im dobre życzenia.

%%K t04-c2-02-1 p
2. --- Mały Piotr jest jeszcze bardzo młody, ma tylko trzy lata. Widzi, jak
\VUgras{kot} \textsuperscript{(1)} się zbliża. Chce pogłaskać jego piękne jedwabiste
\VUgras{futro} i jego długi \VUgras{ogon}, nie zastanawia się, że pod wdzięcznymi
\VUgras{łapami} kryją się złośliwe ostre pazury, które może go podrapią.

%%K t04-c2-03-1 p
3. --- Jego siostra Gabriela, która podaje mu rękę, jest o wiele poważniejsza, a Marceli
ze swoim wielkim \VUgras{płaszczem} \textsuperscript{(4)} i skórzanym \VUgras{pasem}
\textsuperscript{(5)} wygląda na trochę zmężniałego, mimo swoich krótkich spodenek,
które pozwalają widzieć jego \VUgras{pończochy} \textsuperscript{(6)}.

%%K t04-c2-04-1 p
4. --- Ich ojciec włożył swoje grube zimowe \VUgras{palto} \textsuperscript{(7)} z
\VUgras{futrzanym kołnierzem} \textsuperscript{(8)}, a w swojej \VUgras{kieszeni}
\textsuperscript{(9)} ma chustkę. Jego żona ma swój filcowy kapelusz przybrany
\VUgras{piórami} \textsuperscript{(10)}, swój \VUgras{żakiet} \textsuperscript{(11)},
swoje \VUgras{boa} \textsuperscript{(12)} i swoją \VUgras{mufkę} \textsuperscript{(13)}.

%%K t04-c2-05-1 p
5. --- Jest bardzo zimno, bo panuje zima. \VUgras{Śnieg} \textsuperscript{(19)} padał
przez cały dzień i dachy domów są całkiem białe. Wraz z \VUgras{nocą} zimno staje się
jeszcze ostrzejsze. Na niebie \VUgras{księżyc} \textsuperscript{(17)} i \VUgras{gwiazdy}
\textsuperscript{(18)} świecą i przenikają \VUgras{ciemność}.

%%K t04-tit-8 sub
\VUsaut{4.2mm}
\VUpk{10.2pt}{40}{Choroba.}
\VUsaut{3.0mm}

%%K t04-c2-06-1 p
6. --- Biedny \VUgras{ślepiec} \textsuperscript{(20)}, który przechodzi przez plac przed
\VUgras{apteką} \textsuperscript{(16)}, prowadzony przez swojego \VUgras{pudla}
\textsuperscript{(21)}, jest bardzo nieszczęśliwy. Justyna, \VUgras{służąca}
\textsuperscript{(14)}, przynosi z kuchni \VUgras{miskę} \textsuperscript{(15)} bardzo
gorących \VUgras{ziółek} hojnie osłodzonych.

%%K t04-c2-07-1 p
7. --- \VUgras{Babka} \textsuperscript{(23)}, która chce poszukać swoich
\VUgras{okularów} \textsuperscript{(26)} na stole, żeby przeczytać \VUgras{receptę}
\textsuperscript{(27)}, którą \VUgras{lekarz} \textsuperscript{(25)} właśnie napisał,
podnosi się ze swojego fotela, wspierając się na swojej \VUgras{lasce}
\textsuperscript{(24)}.

%%K t04-c2-08-1 p
8. --- Często cierpi na \VUgras{niedomagania, zawroty głowy, katary, bóle zębów}, jej
nogi są słabe, a jej oczy nie są już bardzo dobre. Mimo to chętnie żartuje ze swojego
starego \VUgras{lekarza}, który zawsze chce jej przepisywać \VUgras{miksturę}
\textsuperscript{(28)}. Dzisiaj jest bardzo radosna, bo ma wokół siebie swoją młodą
rodzinę.

%%K t04-c2-09-1 p
9. --- Jest wygodnie w dobrze zamkniętym pokoju. W marmurowym \VUgras{kominku}
\textsuperscript{(29)} płonie \VUgras{piękny ogień} \textsuperscript{(30)} i czujemy się
szczęśliwi w łagodnym \VUgras{świetle} \VUgras{lampy} \textsuperscript{(31)}, którą
dopiero co zapalono.

%%K t04-tit-9 sub
\VUsaut{4.2mm}
\VUpk{10.2pt}{40}{Dziadek.}
\VUsaut{3.0mm}

%%K t04-c2-10-1 p
10. --- Nawet \VUgras{dziadek} \textsuperscript{(33)} zapomniał, że był chory, że jego
\VUgras{reumatyzm} przykuwa go do jego \VUgras{fotela} \textsuperscript{(34)} i że
jeszcze wczoraj miał \VUgras{gorączkę i omdlenia}. Nie pamięta już, że zeszłej zimy
cierpiał na \VUgras{zapalenie płuc} i że chrapliwy i męczący \VUgras{kaszel} bardzo mu
dokuczał.

%%K t04-c2-10-2 p
A jednak z wielkim trudem go wyleczono. Odtąd czuje się nieco lepiej. Ale jego noga,
którą opiera na \VUgras{poduszce} \textsuperscript{(38)} położonej na małym
\VUgras{taborecie} \textsuperscript{(39)}, także go boli.

%%K t04-c2-11-1 p
11. --- Kiedy jest starannie owinięty w swój ciepły \VUgras{szlafrok}
\textsuperscript{(35)} z grubymi \VUgras{guzikami} \textsuperscript{(36)} z masy
perłowej i kiedy ma obok siebie, żeby mu czytać gazetę, małą \VUgras{sierotkę}
\textsuperscript{(40)} ubraną w biały \VUgras{czepek}, w niebieską \VUgras{sukienkę}
\textsuperscript{(42)} i \VUgras{pelerynę} \textsuperscript{(41)}, nie narzeka.

%%K t04-c2-12-1 p
12. --- Co roku, kiedy \VUgras{kalendarz} \textsuperscript{(43)} znów wskazuje
\VUgras{datę} pierwszego \VUgras{grudnia}, niecierpliwie czeka na swoje
\VUgras{wnuki}, bo wie, że nie zapomną tej ważnej \VUgras{daty}.

%%K t04-c2-13-1 p
13. --- Ze wzruszeniem wspomina czasy bardzo dawne, kiedy sam przychodził złożyć
życzenia sławnemu \VUgras{marszałkowi} cesarskiemu, którego \VUgras{portret}
\textsuperscript{(44)} wisi na ścianie i który jest \VUgras{pradziadkiem} jego dzieci.
\VUsaut{6.0mm}
\VUfilet{22mm}
